% Some benchmarks, such as ARC-AGI \cite{cholletARC}, provide only a very small data regime with highly diverse tasks \cite{conceptarc2023}, making it difficult to achieve sound scientific progress on abilities that are crucial to the challenge yet lacking in modern deep learning approaches.
Our motivation stemmed from attempting to make progress at the original ARC challenge \cite{cholletARC}. Similar to the authors of conceptARC \cite{conceptarc2023}, we observed that the data regime of the ARC-AGI was too small (low data availability), too diverse (high variance between individual tasks) to make sound scientific progress on some of the abilities crucial to the challenge and lacking from modern deep learning approaches.
While other researchers adopted the path of scale, big-data regime and novel inference-time methods such as test-time-training and chain of thoughts reasoning to tackle the challenge \cite{Greenblatt_2024,cole_osman_deep_learning_for_arc,OpenAI_2024}, we chose instead to follow a more principled approach, as this allows us to directly evaluate the capabilities of raw architectures and inductive biases to compose and generalize on basic sets of problems. As such, we designed a generator tailored to gauge the systematic and compositional generalization abilities of models, specifically through composing object-centric transformations. Our primary tasks consist of rule-based input-output pairs, where each rule is defined by applying a sequence of transformations to objects arranged on grids of variable sizes.  This setup enables the creation of a wide spectrum of tasks with a large difficulty range. Our framework becomes valuable when varying the transformation sequences and environment parametrizations between training and testing sets -- which
is how we propose to gauge models’ compositional. Specifically, we control systematic and compositional generalization along two axes:

\begin{itemize}[noitemsep, topsep=0pt, leftmargin=1em]
    \item \textbf{Compositional Generalization}: \generatorName\ enables the composition of multiple object transformations at various levels of depth (i.e., the number of transformations applied sequentially) through a set of \ntransforms\ primitive transformations (e.g., translations, rotations, mirroring).
    \item \textbf{Environmental Generalization}: \generatorName\ allows control over additional parameters to assess the generalization capabilities of models, i.e., their ability to apply learned transformations in varied environments. Users can modify the number of objects per grid, the complexity, size, and color of shapes, and the grid size, enabling the generation of different environment versions between training and testing phases.
\end{itemize}

Both axes of research are also enabled through our RGB rendering and sequential rendering, which are respectively introduced in Appendix \ref{appendix:rgb_cogitao} and \ref{appendix:sequential_COGITAO}. While we view these two extensions as essential for future development in natural vision and World-Model research, the main manuscript focuses primarily on the simpler, grid-based input-transformation-output framework.